\begin{solapartat}
\figura[0.4]{sol_fig1.png}

\punts{0,5} $\vec{F} = I(\vec{L}\times\vec{B}) \;\Rightarrow\; |\vec{F}| = I\cdot L\cdot |B|\sin(90\text{°}) = 200\cdot 100\cdot 35\times 10^{-6} = 0,70\ \mathrm{N}$

\punts{0,75} La força és perpendicular al camp i al cable, es troba en el pla de la figura formant un angle de 50° respecte la vertical i sentit cap avall. Per justificar el resultat es pot fer el producte vectorial a partir de les components o bé es pot argumentar que s'ha emprat la regla de la mà dreta.
\end{solapartat}

\begin{solapartat}
\punts{0,4} En les espires A i B el camp magnètic és tangent a la superfície (l'angle entre $\vec{B}$ i $\vec{S}$ és 90° o 270°), de manera que el flux és zero: $\phi = \vec{B}\cdot\vec{S} = B\cdot S\cdot\cos(90\text{°}) = 0\ \mathrm{Wb}$.

Com que el flux és nul, llavors no hi haurà cap corrent induït.

\punts{0,4} A l'espira C el camp magnètic és perpendicular a la superfície (angle entre $\vec{B}$ i $\vec{S}$ de 0° o 180°), per tant, el flux no és nul:
\[ \phi = \vec{B}\cdot\vec{S} = B\cdot S\cdot\cos(0\text{°}) = B\cdot S \]

\punts{0,1} El camp magnètic generat pel fil és proporcional a la intensitat de corrent, per tant, com que el corrent oscil·la de forma periòdica, $I = I_\text{màx}\cdot\cos(\omega t)$, llavors el camp magnètic i el flux també oscil·laran de forma periòdica:
\[ \phi = \phi_\text{màx}\cdot\cos(\omega t) \]

\punts{0,2} Segons la llei de Faraday, si el flux que travessa l'espira varia en el temps, llavors s'induirà una força electromotriu:
\[ \varepsilon = -\frac{d\phi}{dt} = \varepsilon_\text{màx}\cdot\sin(\omega t) \]

\punts{0,15} I finalment, segons la llei d'Ohm, si apareix una fem induïda, llavors es genera un corrent elèctric induït
\[ I = \frac{\varepsilon}{R} = I_\text{màx}\cdot\sin(\omega t) \]

No cal que en la resposta es donin les expressions del flux, la fem i el corrent en funció del temps, ni que es mostri que el corrent generat és un corrent altern. El que és rellevant és que s'identifiqui el mecanisme responsable, que s'indiqui que com que el camp elèctric varia en el temps, llavors, segons la llei de Faraday, s'induirà una fem, i, finalment, que segons la llei d'Ohm, això resultarà en el fet que circularà un corrent elèctric a través de l'espira.
\end{solapartat}
